% Section 2.8, from lectures/L02/README.md.

\section{Summary}
\label{c2:sec:review}

\needspace{6\baselineskip}
\subsection*{What you should be able to do now}
\begin{itemize}
  \item Derive a minimised equation from a Karnaugh map and turn it into a VHDL architecture.
  \item Read and write ports and signals of type \code{std_logic_vector}.
  \item Write a \code{process} with a correct sensitivity list and a \code{case} that covers every
    input value.
  \item Build a design out of more than one \code{entity} with a positional \code{port map}, and
    say why two instances of the same module put two copies of its gates on the FPGA.
  \item Verify a module by running its testbench, and read a failed \code{assert}.
\end{itemize}

\needspace{6\baselineskip}
\subsection*{Questions to test yourself with}
\begin{enumerate}
  \item Given a truth table, can you derive its equation both directly and by way of a Karnaugh
    map? What is the difference between the two results?
  \item Why can a \code{case} not stand directly in an architecture body?
  \item What does \code{when others} protect you from, when the selector ``can only'' be \code{0}
    or \code{1}?
  \item In \code{display1: entity work.display}, what does \code{display1} name, and what does
    \code{work.display} name? Could a second instance reuse the same label?
  \item \code{hex_display}'s architecture contains no logic at all. Where do the gates in the
    finished design come from?
  \item You instantiated \code{display} twice instead of writing its \code{case} twice. Does the
    FPGA get one copy of that logic or two? So what did you actually gain?
  \item The book always writes \code{port map} positionally. What breaks if an entity's ports are
    declared in a different order from the one the testbench expects, and when do you notice?
\end{enumerate}

\needspace{6\baselineskip}
\subsection*{Looking ahead}
\Chapref{c3:ch} goes on with:
\begin{itemize}
  \item Sequential networks: circuits that remember.
  \item The D latch, and the D flip-flop built from a pair of them.
  \item Clocking, and when a signal assignment actually takes effect.
  \item Edge detection: a button press that toggles an LED.
\end{itemize}
